\begin{abstract}
Reinforcement learning with verifiable rewards (RLVR) improves
mathematical reasoning in large language models by rewarding correct
final answers, yet it treats intermediate reasoning as an opaque
black box. This outcome-only supervision leaves reasoning structure
entirely unconstrained, leading to chains that reach correct answers
through unsound derivations, circular logic, or excessive verbosity.
We propose \textbf{TopoPRM}, a verifiable process reward framework
that adds two topology-aware signals to the standard outcome reward.
A rule-based parser extracts the implicit directed acyclic graph (DAG)
of inter-step dependencies from each reasoning trace. A
\emph{topological structure reward} then checks whether this DAG
satisfies acyclicity, connectivity, and directional consistency, while
a \emph{continuity reward} verifies that each step references
expressions or conclusions established earlier. Both signals are
computed by deterministic programs and require no learned reward model
or human annotations. To prevent models from accumulating process reward
at the expense of answer correctness, we introduce Stratified Clipping
Advantage Estimation (SCAE), which partitions completions by
correctness before computing advantages. We train Qwen3-32B with
LoRA-based GRPO using the MS-Swift framework on 637 mathematical
critique prompts and evaluate on 7{,}595 test problems spanning middle
and high school difficulty. TopoPRM produces reasoning traces that
are shorter, more structurally coherent, and more accurate than
those from outcome-only GRPO. A reverse-KL distillation stage transfers
these gains to smaller student models at reduced inference cost.
\end{abstract}
